\documentclass[12pt,a4paper]{article}

%=================================================
% PACKAGES
%=================================================
\usepackage[a4paper, top=2.5cm, bottom=2.5cm, left=2.5cm, right=2.5cm, headheight=15pt]{geometry}
\usepackage{graphicx}
\usepackage{amsmath,amssymb}
\usepackage{float}
\usepackage{xcolor}
\usepackage{fancyhdr}
\usepackage{titlesec}
\usepackage{listings}
\usepackage{hyperref}
\usepackage{tcolorbox}
\usepackage{booktabs}
\usepackage{array}
\usepackage{enumitem}
\usepackage{mdframed}
\usepackage{multicol}
\usepackage{tikz}
\usepackage{circuitikz}
\usepackage[T1]{fontenc}
\usepackage[utf8]{inputenc}
\usepackage{colortbl}

\tcbuselibrary{skins, breakable}

%=================================================
% REUSABLE CIRCUIT DIAGRAM COMMAND
% Usage: \xnorcircuit{scale}
%=================================================
\newcommand{\xnorcircuit}[1]{%
\begin{circuitikz}[scale=#1, transform shape]
\ctikzset{logic ports=ieee}
\def\xa{0}
\def\xdotb{1.2}
\def\xdota{2.2}
\def\xnot{3.5}
\def\xandtwo{5.0}
\def\xandone{7.5}
\def\xor{10.5}

\node[not port] (not1) at (\xnot, 4) {};
\node[not port] (not2) at (\xnot, 2) {};
\node[and port] (and1) at (\xandone, 3) {};
\node[and port] (and2) at (\xandtwo, 0) {};
\node[or port, anchor=in 1] (or1) at (\xor, 3) {};

\draw (\xa,4) node[ocirc]{} node[left=2pt]{$A$}
      -- (\xdota,4) node[circ](dotA){}
      -- (not1.in);
\draw (\xa,2) node[ocirc]{} node[left=2pt]{$B$}
      -- (\xdotb,2) node[circ](dotB){}
      -- (not2.in);

\draw (dotA) |- (and2.in 1);
\draw (dotB) |- (and2.in 2);

\draw (not1.out) -- ++(0.5,0) |- (and1.in 1);
\draw (not2.out) -- ++(0.5,0) |- (and1.in 2);

\draw (and1.out) -- (or1.in 1);
\draw (and2.out) -- (\xor-1.0,0) |- (or1.in 2);

\draw (or1.out) -- ++(0.8,0) node[ocirc]{} node[right=2pt]{$X$};
\end{circuitikz}%
}

%=================================================
% COLORS  (matching reference PDF palette)
%=================================================
\definecolor{accentblue}{RGB}{0, 86, 179}
\definecolor{sectiongray}{RGB}{40, 40, 40}
\definecolor{tableheadbg}{RGB}{0, 86, 179}
\definecolor{rowaltbg}{RGB}{235, 241, 251}
\definecolor{abstractbg}{RGB}{240, 245, 255}
\definecolor{checkgreen}{RGB}{0, 140, 60}
\definecolor{codebg}{RGB}{248, 248, 248}
\definecolor{coderule}{RGB}{180, 180, 180}
\definecolor{titledark}{RGB}{20, 20, 20}

%=================================================
% HYPERREF
%=================================================
\hypersetup{
    colorlinks=true,
    linkcolor=accentblue,
    urlcolor=accentblue,
    pdftitle={GATE IN 2017 -- XNOR Logic Circuit},
    pdfauthor={Shreya Agarwal}
}

%=================================================
% HEADER & FOOTER
%=================================================
\pagestyle{fancy}
\fancyhf{}
\renewcommand{\headrulewidth}{0.4pt}
\renewcommand{\footrulewidth}{0.4pt}
\fancyhead[L]{\small\textcolor{accentblue}{\textbf{GATE IN 2017} Logic Circuit Analysis}}
\fancyhead[R]{\small\textcolor{accentblue}{Digital Electronics}}
\fancyfoot[C]{\small\thepage}

%=================================================
% SECTION STYLE  (matches reference: ALL-CAPS, ruled)
%=================================================
\titleformat{\section}
    {\large\bfseries\color{sectiongray}}
    {\thesection.}{0.6em}{\MakeUppercase}
    [\vspace{-4pt}\textcolor{accentblue}{\rule{\linewidth}{1.2pt}}]

\titlespacing{\section}{0pt}{18pt}{8pt}

\titleformat{\subsection}
    {\normalsize\bfseries\color{accentblue}}
    {\thesubsection}{0.6em}{}

\titlespacing{\subsection}{0pt}{12pt}{4pt}

%=================================================
% CODE / LISTING STYLE
%=================================================
\lstdefinestyle{arduinostyle}{
    language=C++,
    basicstyle=\ttfamily\footnotesize,
    keywordstyle=\color{accentblue}\bfseries,
    commentstyle=\color{gray}\itshape,
    stringstyle=\color{checkgreen},
    numbers=left,
    numberstyle=\tiny\color{gray},
    numbersep=8pt,
    backgroundcolor=\color{codebg},
    frame=single,
    rulecolor=\color{coderule},
    breaklines=true,
    tabsize=4,
    showstringspaces=false,
    captionpos=b
}

\lstdefinestyle{asmstyle}{
    basicstyle=\ttfamily\footnotesize,
    keywordstyle=\color{accentblue}\bfseries,
    commentstyle=\color{gray}\itshape,
    numbers=left,
    numberstyle=\tiny\color{gray},
    numbersep=8pt,
    backgroundcolor=\color{codebg},
    frame=single,
    rulecolor=\color{coderule},
    breaklines=true,
    tabsize=4,
    showstringspaces=false,
    captionpos=b
}

%=================================================
% CUSTOM TABLE STYLE HELPERS
%=================================================
\newcolumntype{C}[1]{>{\centering\arraybackslash}p{#1}}
\newcolumntype{L}[1]{>{\raggedright\arraybackslash}p{#1}}

%=================================================
% TCOLORBOX STYLES
%=================================================
\tcbset{
    abstractbox/.style={
        colback=abstractbg,
        colframe=accentblue,
        fonttitle=\bfseries\color{white},
        coltitle=white,
        attach boxed title to top left={yshift=-2mm, xshift=4mm},
        boxed title style={colback=accentblue, colframe=accentblue},
        boxrule=0.8pt,
        arc=4pt,
        left=8pt, right=8pt, top=8pt, bottom=8pt
    },
    resultbox/.style={
        colback=white,
        colframe=checkgreen,
        fonttitle=\bfseries\color{white},
        coltitle=white,
        attach boxed title to top left={yshift=-2mm, xshift=4mm},
        boxed title style={colback=checkgreen},
        boxrule=1pt,
        arc=4pt,
        left=8pt, right=8pt, top=8pt, bottom=8pt
    }
}

%=================================================
% DOCUMENT
%=================================================
\begin{document}
\thispagestyle{empty}

%=================================================
% COVER PAGE
%=================================================
\begin{titlepage}
\begin{center}

% --- TOP HEADER BAR ---
\colorbox{accentblue}{\parbox{\linewidth}{%
  \centering\vspace{6pt}
  {\large\bfseries\color{white} DEPARTMENT OF ELECTRONICS \& COMMUNICATION ENGINEERING}\\[3pt]
  {\normalsize\color{white} B.E.\ $\cdot$ Academic Year 2026--27}
  \vspace{6pt}
}}

\vspace{0.6cm}

% --- INSTITUTE NAME ---
{\large\bfseries\color{sectiongray} National Institute of Technology Bhopal}

\vspace{0.5cm}
\textcolor{accentblue}{\rule{0.9\linewidth}{1.2pt}}
\vspace{0.4cm}

% --- STAR SYMBOL (matching reference) ---
{\Large\textcolor{accentblue}{$\bigstar$}}

\vspace{0.3cm}

% --- MAIN TITLE ---
{\LARGE\bfseries\color{titledark}
Analysis, Simplification and\\[4pt]
Arduino-Based Verification\\[4pt]
of XNOR Logic Circuit
}

\vspace{0.3cm}

{\large\textcolor{accentblue}{GATE IN 2017 -- Digital Logic Problem}}

\vspace{0.4cm}
\textcolor{accentblue}{\rule{0.9\linewidth}{1.2pt}}

\vspace{0.5cm}

% --- CIRCUIT DIAGRAM ---
\xnorcircuit{0.72}

\vspace{0.5cm}

% --- SUBMISSION DETAILS BOX ---
\begin{tcolorbox}[
  colback=white, colframe=accentblue,
  boxrule=0.8pt, arc=3pt,
  left=12pt, right=12pt, top=6pt, bottom=6pt,
  width=0.72\linewidth
]
\centering
{\bfseries\color{accentblue}\MakeUppercase{Documentation}}\\[6pt]
\begin{tabular}{@{}r@{\quad}l@{}}
\textbf{Submitted by} & Shreya Agarwal \\[3pt]
\textbf{Subject}      & Digital Electronics \& Logic Design \\[3pt]
\textbf{Date}         & \today
\end{tabular}
\end{tcolorbox}

\vfill


\end{center}
\end{titlepage}

%=================================================
% TABLE OF CONTENTS
%=================================================
\tableofcontents
\newpage

%=================================================
\section{Abstract}

\begin{tcolorbox}[abstractbox, title={Overview}]
This report presents a systematic \textbf{analysis, simplification, and hardware
verification} of a digital logic circuit drawn from the GATE Instrumentation
Engineering (IN) 2017 examination paper. The circuit comprises two NOT gates,
two AND gates, and one OR gate. Boolean algebra is applied to reduce the
network, yielding the result:
\[
\boxed{X = AB + \overline{A}\,\overline{B}}
\]
This expression represents the \textbf{Exclusive-NOR (XNOR)} function, which
evaluates to logic HIGH whenever both inputs are equal.

An \textbf{Arduino UNO} platform is subsequently employed to verify this result
experimentally. An LED connected to pin D13 provides visual confirmation:
the LED glows for input combinations $(0,0)$ and $(1,1)$, and extinguishes for
$(0,1)$ and $(1,0)$, in complete agreement with the theoretical derivation.
\end{tcolorbox}

%=================================================
\section{Objectives}

The following objectives guided the course of this project:

\begin{enumerate}[label=\textcolor{accentblue}{\textbf{\arabic*.}}, leftmargin=*, itemsep=4pt]
    \item Analyse the structure of the given GATE IN 2017 logic circuit.
    \item Derive the corresponding Boolean expression from the schematic.
    \item Simplify the expression and identify the equivalent logic function.
    \item Construct a complete truth table to verify the result.
    \item Implement the logic on an Arduino UNO platform and observe output behaviour.
    \item Write an equivalent AVR Assembly program and verify consistency.
\end{enumerate}

%=================================================
\section{Problem Statement}

Determine the Boolean expression corresponding to the given GATE IN 2017
logic circuit (Fig.\ 3.16) and identify the relationship between inputs $A$, $B$
and output $X$. Verify the result using a truth table, an Arduino C++ program,
and an AVR Assembly implementation.

\begin{figure}[H]
\centering
\xnorcircuit{1.0}
\caption{Given GATE IN 2017 Logic Circuit (Fig.\ 3.16)}
\label{fig:circuit}
\end{figure}

\noindent
The multiple-choice options given in the examination are:
\begin{enumerate}[label=\alph*), leftmargin=2cm]
    \item $X = \overline{A}B + \overline{B}A$
    \item $X = AB + \overline{B}A$
    \item $X = AB + \overline{B}A$
    \item $X = \overline{A}B + \overline{B}A$
\end{enumerate}

%=================================================
\section{Circuit Analysis and Simplification}

\subsection{Step-by-Step Derivation}

\noindent\textbf{Step 1 --- NOT Gate Outputs}

The two inverters produce the complements of the inputs:
\[
\overline{A} \qquad \text{and} \qquad \overline{B}
\]

\noindent\textbf{Step 2 --- Upper AND Gate Output}

The AND gate receiving both complements gives:
\[
X_1 = \overline{A}\,\overline{B}
\]

\noindent\textbf{Step 3 --- Lower AND Gate Output}

The AND gate receiving the original inputs gives:
\[
X_2 = AB
\]

\noindent\textbf{Step 4 --- Final OR Gate Output}

The OR gate combines $X_1$ and $X_2$:
\[
X = X_1 + X_2 = \overline{A}\,\overline{B} + AB
\]

\begin{tcolorbox}[resultbox, title={Verified Result}]
\[
\boxed{X = AB + \overline{A}\,\overline{B}}
\]
The circuit implements an \textbf{XNOR} gate. The output is HIGH when both
inputs are equal ($A = B$), and LOW when they differ ($A \neq B$).
\end{tcolorbox}

%=================================================
\section{Truth Table Verification}

\begin{table}[H]
\centering
\caption{Complete Truth Table --- XNOR Function}
\label{tab:truth}
\renewcommand{\arraystretch}{1.4}
\begin{tabular}{
    C{2cm} C{2cm} C{2cm} C{3.5cm}
}
\arrayrulecolor{accentblue}
\toprule
\rowcolor{tableheadbg}
\textcolor{white}{\textbf{A}} &
\textcolor{white}{\textbf{B}} &
\textcolor{white}{\textbf{X}} &
\textcolor{white}{\textbf{Remark}} \\
\midrule
0 & 0 & \textcolor{checkgreen}{\textbf{1}} & \textcolor{checkgreen}{$A = B$ \checkmark} \\
\rowcolor{rowaltbg}
0 & 1 & 0 & $A \neq B$ \\
1 & 0 & 0 & $A \neq B$ \\
\rowcolor{rowaltbg}
1 & 1 & \textcolor{checkgreen}{\textbf{1}} & \textcolor{checkgreen}{$A = B$ \checkmark} \\
\bottomrule
\end{tabular}
\end{table}

\noindent
The truth table confirms that $X = 1$ if and only if $A = B$, which is the defining
property of the XNOR function.

%=================================================
\section{Hardware Implementation}

\subsection{Components Required}

\begin{table}[H]
\centering
\caption{Bill of Materials}
\renewcommand{\arraystretch}{1.4}
\begin{tabular}{L{6cm} C{3cm}}
\arrayrulecolor{accentblue}
\toprule
\rowcolor{tableheadbg}
\textcolor{white}{\textbf{Component}} &
\textcolor{white}{\textbf{Quantity}} \\
\midrule
Arduino UNO Microcontroller Board & 1 \\
\rowcolor{rowaltbg}
Light-Emitting Diode (LED)         & 1 \\
220\,$\Omega$ Current-Limiting Resistor & 1 \\
\rowcolor{rowaltbg}
Breadboard                         & 1 \\
Jumper Wires                       & As required \\
\bottomrule
\end{tabular}
\end{table}

\subsection{Pin Assignment}

\begin{table}[H]
\centering
\caption{Arduino UNO Pin Assignment}
\renewcommand{\arraystretch}{1.4}
\begin{tabular}{C{2.5cm} C{2.5cm} L{6.5cm}}
\arrayrulecolor{accentblue}
\toprule
\rowcolor{tableheadbg}
\textcolor{white}{\textbf{Arduino Pin}} &
\textcolor{white}{\textbf{Signal}} &
\textcolor{white}{\textbf{Function}} \\
\midrule
\textcolor{accentblue}{\textbf{D11}} & Input $A$ LED  & LED indicator for input variable $A$ \\
\rowcolor{rowaltbg}
\textcolor{accentblue}{\textbf{D12}} & Input $B$ LED  & LED indicator for input variable $B$ \\
\textcolor{accentblue}{\textbf{D13}} & Output $X$ LED & LED indicator --- HIGH when $A = B$ (XNOR) \\
\rowcolor{rowaltbg}
\textcolor{accentblue}{\textbf{5V}}  & VCC            & Power supply to LEDs and resistors \\
\textcolor{accentblue}{\textbf{GND}} & GND            & Common ground reference \\
\bottomrule
\end{tabular}
\end{table}

\subsection{Circuit Schematic}

\begin{figure}[H]
\centering
\begin{circuitikz}[scale=0.9, transform shape, line width=0.6pt]

%----------------------------------------------------------
% ARDUINO UNO box  (x: 0..3,  y: 0..9)
%----------------------------------------------------------
\draw (0,0) rectangle (3,9);
\node[font=\bfseries\small] at (1.5,8.5) {Arduino UNO};

% D11  y=7.0
\node[right, font=\small] at (0.2,7.0) {D11};
\draw (3,7.0) -- (3.8,7.0);

% D12  y=5.0
\node[right, font=\small] at (0.2,5.0) {D12};
\draw (3,5.0) -- (3.8,5.0);

% D13  y=3.0
\node[right, font=\small] at (0.2,3.0) {D13};
\draw (3,3.0) -- (3.8,3.0);

% GND  y=1.0
\node[right, font=\small] at (0.2,1.0) {GND};
\draw (3,1.0) -- (3.8,1.0);

%----------------------------------------------------------
% RESISTORS  (3.8 → 6.8)
%----------------------------------------------------------
\draw (3.8,7.0) to[R, l=$R_1{=}220\,\Omega$] (6.8,7.0);
\draw (3.8,5.0) to[R, l=$R_2{=}220\,\Omega$] (6.8,5.0);
\draw (3.8,3.0) to[R, l=$R_3{=}220\,\Omega$] (6.8,3.0);

%----------------------------------------------------------
% LEDs  (6.8 → 9.2)
%----------------------------------------------------------
\draw (6.8,7.0) to[leDo, color=black] (9.2,7.0);
\draw (6.8,5.0) to[leDo, color=black] (9.2,5.0);
\draw (6.8,3.0) to[leDo, color=black] (9.2,3.0);

% LED name labels — placed BELOW each row, clear of symbols
\node[below, font=\small\itshape] at (8.0,6.85) {LED$_A$ — Input $A$};
\node[below, font=\small\itshape] at (8.0,4.85) {LED$_B$ — Input $B$};
\node[below, font=\small\itshape] at (8.0,2.85) {LED$_X$ — Output $X$};

%----------------------------------------------------------
% GND bus  at y = -0.5
%----------------------------------------------------------
% LED cathodes drop straight down to GND bus
\draw (9.2,7.0) -- (9.2,-0.5);
\draw (9.2,5.0) -- (9.2,-0.5);   % junction
\draw (9.2,3.0) -- (9.2,-0.5);   % junction

% Horizontal GND bus
\draw (3.8,-0.5) -- (9.2,-0.5);

% Arduino GND drops to bus
\draw (3.8,1.0) -- (3.8,-0.5);

% Ground symbol
\draw (9.2,-0.5) node[ground]{};

%----------------------------------------------------------
% JUNCTION DOTS
%----------------------------------------------------------
\filldraw (9.2,5.0)  circle (2pt);
\filldraw (9.2,3.0)  circle (2pt);
\filldraw (3.8,-0.5) circle (2pt);

\end{circuitikz}
\caption{Arduino UNO Implementation Schematic}
\label{fig:schematic}
\end{figure}

\subsection{Working Principle}

The Arduino firmware iterates through all four possible input combinations of
$A$ and $B$ in a loop. LEDs on pins D11 and D12 display the current values of
inputs $A$ and $B$ respectively, so the input state is always visible. The
XNOR logic is computed in software and the result is driven on pin D13. The
output LED glows whenever $X = 1$ (i.e.\ $A = B$) and extinguishes whenever
$X = 0$ (i.e.\ $A \neq B$).

%=================================================
\newpage
\section{Arduino Program}

\begin{lstlisting}[style=arduinostyle, caption={XNOR Verification --- Arduino C++ Program}]
const int A_LED = 11;   // Input A  -- LED indicator
const int B_LED = 12;   // Input B  -- LED indicator
const int Y_LED = 13;   // Output X -- LED indicator

void setup()
{
    pinMode(A_LED, OUTPUT);
    pinMode(B_LED, OUTPUT);
    pinMode(Y_LED, OUTPUT);
    Serial.begin(9600);
    Serial.println("A  B  X (XNOR)");
}

void loop()
{
    for (int i = 0; i < 4; i++)
    {
        int A = (i >> 1) & 1;  // Extract bit 1
        int B =  i       & 1;  // Extract bit 0

        // XNOR: output HIGH when A == B
        int X = (A == B) ? 1 : 0;

        digitalWrite(A_LED, A);   // Show current A on LED
        digitalWrite(B_LED, B);   // Show current B on LED
        digitalWrite(Y_LED, X);   // Show XNOR output on LED

        Serial.print(A); Serial.print("  ");
        Serial.print(B); Serial.print("  ");
        Serial.println(X);

        delay(1000);            // Hold for 1 second
    }
}
\end{lstlisting}

%=================================================
\newpage
\section{AVR Assembly Program}

\begin{lstlisting}[style=asmstyle, caption={XNOR Verification --- AVR Assembly (ATmega328P)}]
; Pin mapping:
;   PB3 = Arduino D11 -> A LED
;   PB4 = Arduino D12 -> B LED
;   PB5 = Arduino D13 -> X LED (XNOR output)

.include "m328pdef.inc"

.org 0x0000
    rjmp MAIN

MAIN:
    ldi  r16, 0b00111000    ; Set PB3, PB4, PB5 as outputs
    out  DDRB, r16

LOOP:
    ldi  r17, 0             ; Loop counter i = 0

NEXT:
    ; --- Extract A (bit1) and B (bit0) from counter ---
    mov  r18, r17
    lsr  r18
    andi r18, 0x01          ; r18 = A

    mov  r19, r17
    andi r19, 0x01          ; r19 = B

    ; --- Drive A LED on PB3 ---
    tst  r18
    breq A_LOW
    sbi  PORTB, 3           ; A = 1 -> LED ON
    rjmp A_DONE
A_LOW:
    cbi  PORTB, 3           ; A = 0 -> LED OFF
A_DONE:

    ; --- Drive B LED on PB4 ---
    tst  r19
    breq B_LOW
    sbi  PORTB, 4           ; B = 1 -> LED ON
    rjmp B_DONE
B_LOW:
    cbi  PORTB, 4           ; B = 0 -> LED OFF
B_DONE:

    ; --- Compute XNOR: X = 1 if A == B ---
    cp   r18, r19
    breq SET_HIGH
    cbi  PORTB, 5           ; A != B -> X LED OFF
    rjmp DONE
SET_HIGH:
    sbi  PORTB, 5           ; A == B -> X LED ON

DONE:
    rcall DELAY             ; ~1 second delay

    inc   r17
    cpi   r17, 4
    brne  NEXT              ; Repeat for all 4 combinations

    rjmp  LOOP              ; Restart cycle indefinitely

DELAY:
    ldi  r20, 100
L1: ldi  r21, 255
L2: ldi  r22, 255
L3: dec  r22
    brne L3
    dec  r21
    brne L2
    dec  r20
    brne L1
    ret
\end{lstlisting}

%=================================================
\section{Results}

\begin{tcolorbox}[abstractbox, title={Experimental Findings}]
The Arduino UNO successfully iterated through all four input combinations of
$A$ and $B$. LEDs on D11 and D12 visually confirmed the current input state,
while the output LED on D13 matched the XNOR truth table exactly:
\begin{itemize}[itemsep=2pt]
    \item \textbf{D11 (A LED)} and \textbf{D12 (B LED)} toggled with each combination \quad
          \textcolor{checkgreen}{\checkmark}
    \item \textbf{D13 (X LED) ON}  for $(A,B) = (0,0)$ and $(1,1)$ \quad
          \textcolor{checkgreen}{\checkmark}
    \item \textbf{D13 (X LED) OFF} for $(A,B) = (0,1)$ and $(1,0)$ \quad
          \textcolor{checkgreen}{\checkmark}
\end{itemize}
Both the Arduino C++ firmware and the AVR Assembly implementation produced
identical results across all three LEDs, demonstrating consistency at both
levels of abstraction.
\end{tcolorbox}

%=================================================
\section{Optimisation Analysis}

\subsection{Gate Count Comparison}

\begin{table}[H]
\centering
\caption{Original Circuit vs.\ Direct XNOR Gate Implementation}
\renewcommand{\arraystretch}{1.4}
\begin{tabular}{L{5.5cm} C{3cm} C{3cm}}
\arrayrulecolor{accentblue}
\toprule
\rowcolor{tableheadbg}
\textcolor{white}{\textbf{Metric}} &
\textcolor{white}{\textbf{Original Circuit}} &
\textcolor{white}{\textbf{Simplified}} \\
\midrule
NOT Gates            & 2 & 0 \\
\rowcolor{rowaltbg}
AND Gates            & 2 & 0 \\
OR Gates             & 1 & 0 \\
\rowcolor{rowaltbg}
Total Gates          & 5 & 1 (XNOR IC) \\
Propagation Delay    & 3 stages & 1 stage \\
\rowcolor{rowaltbg}
Implementation       & Discrete gates & 74HCXX / CMOS \\
\bottomrule
\end{tabular}
\end{table}

\subsection{Advantages of Simplification}

\begin{itemize}[leftmargin=*, itemsep=4pt,
    label=\textcolor{accentblue}{$\blacktriangleright$}]
    \item \textbf{Reduced gate count} --- five discrete gates replaced by a
          single XNOR gate (e.g.\ 74HC266), lowering component cost.
    \item \textbf{Lower propagation delay} --- three cascaded gate delays
          reduced to one, improving switching speed.
    \item \textbf{Reduced power consumption} --- fewer active logic elements
          draw less quiescent current.
    \item \textbf{Improved reliability} --- fewer components result in fewer
          potential failure points and higher MTBF.
    \item \textbf{Smaller PCB footprint} --- significant area savings in
          integrated circuit layouts and printed circuit board designs.
\end{itemize}

%=================================================
\section{Conclusion}

The given GATE IN 2017 logic circuit was systematically analysed using Boolean
algebra. The step-by-step derivation established that:
\[
X = AB + \overline{A}\,\overline{B}
\]
This is the \textbf{XNOR} function --- the output is logic HIGH when both inputs
are equal, and logic LOW when they differ.

The truth table, Arduino C++ firmware, and AVR Assembly program all confirmed
this result experimentally. This project underscores the practical value of
Boolean simplification and gate-level optimisation in reducing hardware cost,
minimising propagation delay, and streamlining design complexity in digital
electronic systems.

\vfill
\begin{center}
\textcolor{accentblue}{\rule{0.6\linewidth}{0.8pt}}\\[4pt]
\textit{--- End of Report ---}
\end{center}

\end{document}





















































